\section{Preliminaries: What Does It Mean for a Qubit to Be ``Clean''?}
\label{sec:preliminaries}

\subsection{Known state, separability, and liveness}

\begin{definition}[Clean qubit]
A qubit $q$ is \emph{clean with value $\ket{\phi}$} relative to the remaining
live system $L$ when the joint state factors as
\begin{equation}
  \rho_{qL}=\ket{\phi}\!\bra{\phi}_{q}\otimes\rho_L.
  \label{eq:clean-definition}
\end{equation}
\end{definition}

The factorization is essential. A qubit that produces $0$ with high probability
is not necessarily clean. It may retain phase error, leakage outside the
computational subspace, or correlations with live data.

\begin{definition}[Reclaimed qubit]
A clean qubit is \emph{reclaimed} when the legitimate program has no future
semantic dependence on its previous logical identity and the runtime assigns it
to a new logical role.
\end{definition}

The distinction mirrors classical memory management, but quantum cleanup is
stricter. A classical buffer can be overwritten even when correlated copies
exist elsewhere. A quantum qubit cannot be safely released if it remains
entangled with live outputs.

\subsection{Reduced states and hidden correlations}

Suppose a legitimate register $S$ shares a joint state $\rho_{SE}$ with an
attacker-controlled register $E$. The legitimate user sees only
\begin{equation}
  \rho_S=\Tr_E(\rho_{SE}).
\end{equation}
Two global states can have the same reduced state on $S$ while differing in
correlations with $E$. This is the mathematical location of a possible hidden
channel.

Yet there is an important constraint: when $\rho_S$ is a pure state, any joint
extension must factor. Thus, if an arbitrary pure victim state is exactly
preserved, an environment cannot remain correlated with it. Hidden leakage is
possible only when the inputs contain an accessible classical component, the
victim checks an incomplete observable, or the implementation exposes
additional preparation information.
